\printedpage{29}
% head: A. Grothendieck : Dérivateurs, Ch. II, Édition des Universités de Montpellier II et Paris VII
% foot: Version 8 juillet 2001

\msp{page 45}
is more natural to take $M = e$ (one-point set) and to interpret the arrow (3.2.3) as
\[
\Hom(e_I, F) \longrightarrow \Hom(e_{I'}, u^*F)
\tag*{(3.2.4)}
\]
and by adjunction of $u_!$, $u^*$ the second member is interpreted as $\Hom(u_!(e_{I'}), F)$, and the arrow is deduced from
\[
u_!e_{I'} = u_!u^*(e_I) \xrightarrow{\text{adj.}} e_I
\tag*{(3.2.5)}
\]
by taking Hom into $F$. Thus to say that (3.2.4) is \texttt{[an]} isomorphism for every $F$ means that (3.2.5) is \texttt{[an]} isomorphism. Passing to the values of the two members at $i$, we find $\pi_0(i\bs I') \xrightarrow{\sim} 1$, i.e. that $i\bs I'$ is 0-connected, q.e.d.

\noindent\textbf{Definition 3.3.}
\begin{itemize}[label={}]
\item[a)] When the equivalent conditions of 3.2 are satisfied, thus when the $i\bs I'$ ($i \in \Ob I$) are 0-connected, we say that $u : I' \longrightarrow I$ is a \textit{cofinal} functor. We say that $u$ is \textit{co-cofinal} if $u^o : I'^o \longrightarrow I^o$ is cofinal, i.e. if the localized categories $I'/i$ of $I'$ over $I$ ($i \in \Ob I$) are 0-connected.
\item[b)] If $I'$ is a \textit{subcategory} of $I$, we say that $I'$ \textit{is cofinal in} $I$ if the inclusion functor is cofinal, \textit{co-cofinal} if this functor is co-cofinal.
\end{itemize}

\noindent\textbf{Corollary 3.4.} (\footnote[55]{\textbf{NB} For a complementary cofinality statement, cf. XVIII prop. 4.9.6.3 (p. 224-225, 228-242). A statement covering both would be needed.})
\begin{itemize}[label={}]
\item[\textit{a)}] \textit{For $u$ to be cofinal, it is} necessary \textit{that the colocalized categories $i\bs I'$ be nonempty, i.e. that we have}

\msp{page 46}
\begin{itemize}[label={}]
\item[\textit{F 0)}] $\pi_0(f)$ \texttt{[rather $\pi_0(u)$]} \textit{is bijective.}
\item[\textit{F 1)}] \textit{For every} $i \in \Ob I$, \texttt{[there]} \textit{exists} \texttt{[an]} $i' \in I'$ \textit{and an arrow} $i \longrightarrow u(i')$ \textit{in} $I$.
\end{itemize}
(\footnote[56]{\textbf{NB} F 1 already implies \texttt{[the]} surjectivity of $\pi_0$, hence its bijectivity if $I'$ is 0-connected - in this case F 0 follows from F 1.}).
\item[\textit{b)}] \textit{If $I'$ is weakly filtering, for $u$ to be cofinal, it is} sufficient \textit{that it satisfy F 1, as well as the condition:}
\begin{itemize}[label={}]
\item[\textit{F 2)}] \textit{For every double arrow} $i \dbl{\beta}{\gamma} u(i')$ \textit{(with} $i \in \Ob I$, $i' \in \Ob I'$\textit{),} \texttt{[there]} \textit{exists an arrow} $i' \xrightarrow{\alpha'} j'$ \textit{in} $I'$\textit{, in such a way that} $u(\alpha')$ \textit{equalizes} $(\beta, \gamma)$\textit{, i.e.} $u(\alpha')\beta = u(\alpha')\gamma$.
\end{itemize}
\textit{and these conditions F 1, F 2 imply that the $i\bs I'$ are weakly filtering, and $I$ as well} (\footnote[57]{If $I'$ is locally weakly filtering (= prefiltering), then for the cofinality of $u$, it suffices that we have F 0, F 1, F 2, and then the $i\bs I'$ are weakly filtering, and $I$ is prefiltering just like $I'$.}).

\printedpage{30}
% head: A. Grothendieck : Dérivateurs, Ch. II, Édition des Universités de Montpellier II et Paris VII
% foot: Version 8 juillet 2001
\item[\textit{c)}] \textit{If $I'$ is filtering, then $u$ is cofinal} if and only if \textit{it satisfies the two conditions F1, F2, or again if and only if the $i\bs I'$ are} filtering \textit{(and not merely} 0-\textit{connected). And this implies that $I'$ is also filtering} (\footnote[58]{If $I'$ is locally filtering, $u$ is cofinal if and only if it satisfies F 0, F 1, F 2, i.e. if and only if the $i\bs I'$ are filtering, and then $I$ is locally filtering.}).
\item[\textit{d)}] \textit{If $I$ is prefiltering (= locally weakly filtering, i.e. PS 1), and $u$} fully faithful, \textit{for $u$ to be cofinal, it is} necessary \textit{and it is} sufficient \textit{that $u$ satisfy condition F 1 above. And then $I'$ is also prefiltering, and the $i\bs I'$ are weakly filtering. If moreover $I$ is filtering, the $i\bs I'$ and $I'$ are as well. (Compare SGA 4, I 8.1.3.)}
\end{itemize}

\msp{page 47}
\textsc{Proof.} a) F 1 is obvious, since 0-connected $\Longrightarrow$ nonempty. F 0 because $W_0$-coaspheric $\Longrightarrow$ $W_0$-equivalence, i.e. $\pi_0(u)$ bijective.

To \textit{prove} b), it suffices to prove that F 1, F 2 (if $I'$ is weakly filtering) imply that the $i\bs I'$ are weakly filtering (a fortiori 0-connected) and moreover that $I$ is as well. We already know that they are nonempty by F 1, remains to see that if $(i', \alpha)$, $(j', \beta)$ are in $i\bs I'$, then they are ``dominated'' by a third object of $i\bs I'$. Since $I'$ is weakly filtering, we find
\begin{tikzcd}[row sep=0pt, column sep=small, baseline=(k.base)]
i' \arrow[dr, "r"] & \\
& |[alias=k]| k' \\
j' \arrow[ur, "s"'] &
\end{tikzcd}, whence $u(r)$, $u(s)$ in $I$ (cf. diagram)
\[
\begin{tikzcd}[row sep=small, column sep=small]
& i' \arrow[rr, "r"] \arrow[dd, dash] & & k' \arrow[r, dotted, "t"] \arrow[dd, dash] & l' \arrow[dd, dash] \\
& & j' \arrow[ur, "s"] \arrow[dd, dash] & & \\
& u(i') \arrow[rr, "u(r)"] & & u(k') \arrow[r, dotted, "u(t)"] & u(l') \\
i \arrow[ur, "\alpha"] \arrow[rr, "\beta"'] & & u(j'), \arrow[ur, "u(s)"'] & &
\end{tikzcd}
\]
applying F 2 to the double arrow $u(r)\alpha, u(s)\beta : i \rightrightarrows u(k')$, we find $k' \xrightarrow{t} l'$ (dotted on \texttt{[the]} diagram) such that $(u(t)u(r))\alpha = (u(t)u(s))\beta$, say $\gamma$, which shows that the object $(l', \gamma : i \longrightarrow u(l'))$ dominates $(i', \alpha)$ and $(j', \beta)$.

The proof that if $I'$ is weakly filtering and $u$ satisfies F 1, then $I$ is weakly filtering, is immediate.

\textit{Let us prove} c). It remains to prove that if $I'$ is \textit{filtering} (not merely weakly filtering), then if $u$ \texttt{[is]} cofinal, then F 2 is satisfied, and moreover the $i\bs I'$ and $I$ are filtering (not merely weakly filtering). But let us note at once that $i\bs I' \longrightarrow I'$ is coétale, hence $I'$ filtering, hence locally filtering implies that $i\bs I'$ is also locally filtering\,; since by hypothesis the $i\bs I'$ are 0-connected, they are therefore filtering. Condition F 2 comes out easily now:

\printedpage{31}
% head: A. Grothendieck : Dérivateurs, Ch. II, Édition des Universités de Montpellier II et Paris VII
% foot: Version 8 juillet 2001
\msp{page 48}
\[
\begin{tikzcd}[row sep=small, column sep=small]
& i' \arrow[rr, "r"] \arrow[dd, dash] & & j' \arrow[r, dotted, "t"] \arrow[dd, dash] & k' \arrow[dd, dash] \\
& & i' \arrow[ur, "s"] \arrow[dd, dash] & & \\
& u(i') \arrow[rr, "u(r)"] & & u(j') \arrow[r, dotted, "u(t)"] & u(k') \\
i \arrow[ur, "\alpha"] \arrow[rr, "\beta"'] & & u(i'). \arrow[ur, "u(s)"'] & &
\end{tikzcd}
\]
Since $i\bs I'$ is filtering, hence weakly filtering (PS 1), we can find $i' \dbl{r}{s} j'$ such that
\[
u(r)\alpha = u(s)\beta,
\]
and since $I'$ is filtering, we can find $t : j' \longrightarrow k'$ which equalizes the double arrow $(r, s)$:
\[
tr = ts(\defeq \tau),
\]
then
\[
\underbrace{u(t)u(r)}_{u(tr)}\alpha = \underbrace{u(t)u(s)}_{u(ts)}\beta \qquad\quad \text{i.e.} \qquad\quad u(\tau)\alpha = u(\tau)\beta,
\]
i.e. we have verified F 2.

Remains to prove that $I$ is filtering (if $I'$ is and \texttt{[if]} $u : I' \longrightarrow I$ is cofinal). We already know that $I$ is weakly filtering by b), remains to verify the double arrow condition
\[
i \dbl{\alpha}{\beta} j
\]
(PS 2). Let $i'$ be in $I'$ and $j \longrightarrow u(i')$, and consider
\[
i \dbl{\gamma\alpha}{\gamma\beta} u(i'),
\]
by hypothesis F 2 $\exists\, r : i' \longrightarrow j'$ in $I'$ such that
\msp{page 49}
$u(r)(\gamma\alpha) = u(r)(\gamma\beta)$, hence $u(r)\gamma$ equalizes $(\alpha, \beta)$, q.e.d.

The verification of d) by AQT \texttt{[abbreviation whose meaning eludes the editors]}.

\printedpage{32}
% head: A. Grothendieck : Dérivateurs, Ch. II, Édition des Universités de Montpellier II et Paris VII
% foot: Version 8 juillet 2001
\section*{4\quad Biaspheric functors}

Let again
\[
u : I' \longrightarrow I
\]
be an arrow of Cat, let $\cM$ be a category, whence
\[
u^* : \cM(I) \longrightarrow \cM(I');
\tag*{(4.1)}
\]
we seek conditions on $u$, $\cM$ for this functor to be fully faithful. If the left adjoint functor $u_!$ (resp. right adjoint $u_*$) exists, we know that it amounts to the same to say that the adjunction morphism
\[
u_!u^*F \longrightarrow F
\tag*{(4.2)}
\]
resp.
\[
u_*u^*F \longleftarrow F
\tag*{(4.2 bis)}
\]
is \texttt{[an]} isomorphism for every $F$ in $\cM(I)$. Let us then say that $u$ is $\cM$-connected, or $\cM$-biaspheric, this terminology implying that this is a property which refers to the derivator $\mathbf{D}_{\cM} = (S \longmapsto \cM(S))$ on Cat defined by $\cM$. Let us note that
\[
\begin{aligned}
&\text{If $u$ is $\cM$-connected, then it is also $\cM(S)$-connected, for every $S$ in Cat,}\\
&\text{since $\cM(S)(I) \simeq \cM(I)(S) \defeq \uHom(S^o, \cM(I))$,}
\end{aligned}
\tag*{(4.3)}
\]

\msp{page 50}
and
\[
\text{If $\cM' \hookrightarrow \cM$ is fully faithful, then $\cM$-connected implies $\cM'$-connected,}
\tag*{(4.4)}
\]
finally, conversely,
\[
\begin{aligned}
&\text{for $u$ to be $\cM$-connected, it is necessary and sufficient that for every \textit{small}}\\
&\text{full subcategory $\cM'$ of $\cM$, $u$ be $\cM'$-connected.}
\end{aligned}
\tag*{(4.5)}
\]
Since every small $\cM'$ embeds in $\cM'^\wedge = \operatorname{Ens}(\cM')$, it follows from (4.3), (4.4), (4.5) that
\[
\text{If $u$ is Ens-connected, it is $\cM$-connected for every category $\cM$,}
\tag*{(4.6)}
\]
for this is the case if $\cM$ is small by (4.3), hence for every $\cM$ by (4.5).

Let us say that a functor $u$ in Cat is $\cM$-\textit{cofinal} or $\cM$-coaspheric if the functorial homomorphism
\[
\varprojlim_{I} F \longrightarrow \varprojlim_{I'} u^*(F)
\tag*{(4.7)}
\]
(arrow in $\uHom(\cM(I), \cM)$ if the $\varprojlim$ of types $I$, $I'$ exist in $\cM$, otherwise arrow in $\uHom(\cM(I), \cM^\wedge)$) is an isomorphism. Reducing to the case of small categories $\cM$, each embedded in $\cM^\wedge$, we see by criterion (4.2 bis) that
\[
\text{$u$ is $\cM$-connected if for every $i \in \Ob I$,}
\tag*{(4.8)}
\]

\printedpage{33}
% head: A. Grothendieck : Dérivateurs, Ch. II, Édition des Universités de Montpellier II et Paris VII
% foot: Version 8 juillet 2001
\msp{page 51}
the induced functor $I'/i \longrightarrow I/i$ is $\cM$-cofinal. The converse is not clear, for $u$ is $\cM$-connected if and only if for every $i \in \Ob I$ and every $\cM$-presheaf $F_i$ on $I/i$ \textit{induced by a presheaf $F$ on $I$,} the homomorphism $\varprojlim\limits_{I/i} F_i \longrightarrow \varprojlim\limits_{I'/i} (u/i)^*(F/i)$ is \texttt{[an]} isomorphism. We can therefore assert straightaway that this implies likewise for \textit{every} $F_i$ on $I/i$, only if we know that $(u/i)^* : \cM(I) \longrightarrow \cM(I/i)$ is essentially surjective. This is so for every $i$ when $I$ is \textit{ordered,} for then $I/i \longrightarrow I$ is fully faithful, hence $\cM(I/i) \longrightarrow \cM(I)$ \texttt{[is]} fully faithful, hence $\cM(I) \longrightarrow \cM(I/i)$ \texttt{[is]} a localization functor, in particular essentially surjective.

\msp{page 52}
But when $\cM = (\text{Ens})$, this perplexity disappears. We find the

\noindent\textbf{Proposition 4.1.} \textit{Let $u : I' \longrightarrow I$ be an arrow in} Cat. \textit{Equivalent conditions:}
\begin{itemize}[label={}]
\item[\textit{a)}] \textit{For every category $\cM$, $u^*_{\cM} : \cM(I) \longrightarrow \cM(I')$ is fully faithful (i.e. $u$ is $\cM$-connected).}
\item[\textit{a$'$)}] \textit{For every category $\cM$ in which direct limits of type $I$, $I'$ exist, the adjunction arrow for $F$ in $\cM(I)$} % Transcription corrected: official printed33 reads a prime.
\[
u_!u^*F \xrightarrow{\sim} F
\]
\textit{is} \texttt{[an]} \textit{isomorphism.}
\item[\textit{a$''$)}] \textit{Dual of a$'$), with} \texttt{[the]} \textit{adjunction arrow}
\[
u_*u^* \xleftarrow{\sim} F. % [sic] u_*u^* -> u_*u^*F, the object F is missing on the left side (compare a') and (4.2 bis))
\]
\item[\textit{b) b$'$) b$''$)}] \textit{As in a) a$'$) a$''$), with $\cM =$} (Ens).
\item[\textit{c)}] \textit{As in b$'$), with $F$ representable, i.e. of the form $j \longmapsto$} $\Hom_I(j, i)$.
\item[\textit{d)}] \textit{For every $I$ in $I$, $I'/i \longrightarrow I/i$ is cofinal.} % [sic] Pour tout I dans I -> Pour tout i dans I, the variable of the statement is i (compare d'))
\item[\textit{d$'$)}] \textit{For every $i$ in $I$, $i\bs I' \longrightarrow i\bs I$ is cofinal.}
\item[\textit{d$''$)}] \textit{For every arrow $i_0 \xrightarrow{\alpha} i_1$ in $I$, the category} $i_0\bs \underset{\alpha}{I'}/i_1 \defeq I' \times_I \underbrace{i_0\bs \underset{\alpha}{I}/i_1}_{\substack{\text{biinduced category}\\ (i_0, \alpha)\bs(I/i_1) \simeq (i_0\bs I)/(i_1, \alpha)}}$ \qquad \textit{is}

0-\textit{connected.}
\end{itemize}

\msp{page 53}
\textsc{Proof.} We already know, from the preliminaries, the equivalence of a, a$'$, a$''$, b, b$'$, b$''$, let A be this condition - we shall say that $u$ is $W_0$-\textit{connected} or $W_0$-\textit{biaspheric,} if these equivalent conditions are satisfied. We have moreover the self-duality
\[
\text{$u$ is $W_0$-connected if and only if $u^o$ is,}
\tag*{(4.8)} % [sic] (4.8) -> another number, the label (4.8) is already used for a display on p. 32
\]
for we obviously have
\[
\text{$u$ is $\cM$-connected if and only if $u^o$ is $\cM^o$-connected.}
\tag*{(4.9)}
\]